\documentclass[CMPE]{KGCOEReport}

\usepackage{graphicx}
\usepackage{newtxtext, newtxmath}
\usepackage{csquotes}
\usepackage{parskip}
\usepackage{karnaugh-map}
\usepackage{float}
\usepackage{hyperref}
\usepackage[super]{nth}
\usepackage{pdfpages}

\graphicspath{{./assets/}}

\newcommand{\classCode}{CMPE 260}
\newcommand{\name}{Aden Perry}
\newcommand{\LabSectionNum}{L2}
\newcommand{\LabInstructor}{Richard Cliver}

\newcommand{\TAs}{Eric Reed \\ Tevin Hendess}
\newcommand{\LectureSectionNum}{01}
\newcommand{\LectureInstructor}{Cory Merkel}
\newcommand{\exerciseNumber}{4}
\newcommand{\exerciseDescription}{Execute Stage}
\newcommand{\dateDone}{2026-03-03}
\newcommand{\dateSubmitted}{2026-03-20}

\begin{document}

\maketitle

\section*{Abstract}
\addcontentsline{toc}{section}{Abstract}

The purpose of this exercise was to design, implement, and test an n-bit
Arithmetic Logic Unit (ALU) to support multiplication, addition, loigcal left-
and right-shifts, arithmetic right-shifts, logical AND, logical OR, and logical
XOR; as well as to create and execute stage to map input flags and register and
memory values to the ALU. The partial n-bit ALU from Laboratory Exercise One
was expanded upon.

\section*{Design Methodology}
\addcontentsline{toc}{section}{Design Methodology}

\section*{Results and Analysis}
\addcontentsline{toc}{section}{Results and Analysis}

\section*{Conclusion}
\addcontentsline{toc}{section}{Conclusion}

\includepdf[pages=-]{assets/prelab}
\end{document}
